\documentclass[10pt,a4paper]{article}
\usepackage[utf8]{inputenc}
\usepackage[vietnamese]{babel}
\usepackage{amsmath,amssymb,amsfonts}
\usepackage{geometry}
\geometry{a4paper, left=10mm, right=10mm, top=10mm, bottom=10mm}
\usepackage{tikz}
\usetikzlibrary{shapes,arrows.meta,positioning,calc}
\usepackage[most]{tcolorbox}
\usepackage{array,booktabs,multicol}
\usepackage{fontawesome}
\usepackage{helvet}
\renewcommand{\familydefault}{\sfdefault}
\usepackage{parskip}

\definecolor{primary}{RGB}{112, 48, 160}
\definecolor{teal}{RGB}{0, 128, 128}
\definecolor{orange}{RGB}{230, 108, 44}
\definecolor{green}{RGB}{46, 139, 87}
\definecolor{rose}{RGB}{192, 0, 0}

\pagestyle{empty}

\begin{document}

% --- HEADER BANNER ---
\begin{tcolorbox}[
  enhanced,
  colback=violet!8!white,
  colframe=violet!80!black,
  arc=2.5mm, boxrule=1pt,
  left=3mm, right=3mm, top=2mm, bottom=2mm
]
  \begin{tabular*}{\textwidth}{@{\extracolsep{\fill}} lr}
    {\small\bfseries\color{violet!85!black}\faBook\ SỔ TAY KIẾN THỨC TOÁN 2} & 
    {\small\bfseries\color{teal!85!black}\faGraduationCap\ CHƯƠNG TRÌNH GDPT 2018}
  \end{tabular*}
  \vspace{-1mm}
  \begin{center}
    \rule{0.98\textwidth}{0.6pt}
  \end{center}
  \vspace{-2mm}
  \centering
  {\LARGE\bfseries\color{violet!90!black}\faBookmark\ BÀI 65: BIỂU ĐỒ TRANH}\\[1mm]
  {\itshape\footnotesize\color{black!75} Kĩ năng đọc, phân tích và biểu diễn số liệu bằng biểu đồ tranh}
\end{tcolorbox}

\vspace{-2mm}

\noindent
\begin{minipage}[t]{0.49\textwidth}\vspace{0pt}
\begin{tcolorbox}[
  enhanced,
  colback=violet!4!white,
  colframe=violet!75!black,
  arc=2mm, boxrule=0.8pt,
  left=2mm, right=2mm, top=1.8mm, bottom=1.8mm,
  title={\footnotesize\bfseries\color{white}\faKey\ 1. Khái niệm biểu đồ tranh}
]
{\footnotesize
\begin{itemize}\setlength{\itemsep}{1pt}\setlength{\leftskip}{-2.5mm}
  \item \textbf{Biểu đồ tranh} là biểu đồ sử dụng hình ảnh, biểu tượng hoặc bức tranh để biểu thị số liệu thu thập được.
  \item Cấu tạo của một biểu đồ tranh gồm:
  \begin{itemize}\setlength{\itemsep}{0.5pt}\setlength{\leftskip}{-2mm}
    \item \textbf{Tên biểu đồ:} Cho biết biểu đồ nói về nội dung gì.
    \item \textbf{Các hàng (hoặc cột):} Tên đối tượng và các biểu tượng tranh tương ứng.
    \item \textbf{Ghi chú:} Mỗi biểu tượng thay cho bao nhiêu đối tượng (ở lớp 2, mỗi biểu tượng thường ứng với \textbf{1 đối tượng}).
  \end{itemize}
\end{itemize}
}
\end{tcolorbox}
\end{minipage}
\hfill
\begin{minipage}[t]{0.49\textwidth}\vspace{0pt}
\begin{tcolorbox}[
  enhanced,
  colback=teal!4!white,
  colframe=teal!75!black,
  arc=2mm, boxrule=0.8pt,
  left=2mm, right=2mm, top=1.8mm, bottom=1.8mm,
  title={\footnotesize\bfseries\color{white}\faSliders\ 2. Các bước đọc và khai thác biểu đồ tranh}
]
{\footnotesize
\begin{itemize}\setlength{\itemsep}{1pt}\setlength{\leftskip}{-2.5mm}
  \item \textbf{Bước 1: Đọc tên biểu đồ} để hiểu chủ đề cần tìm hiểu.
  \item \textbf{Bước 2: Xem dòng quy ước/ghi chú} xem mỗi hình vẽ biểu diễn bao nhiêu đơn vị.
  \item \textbf{Bước 3: Đếm số biểu tượng} ở mỗi dòng để biết số lượng của từng loại đối tượng.
  \item \textbf{Bước 4: So sánh và tính toán:}
  \begin{itemize}\setlength{\itemsep}{0.5pt}\setlength{\leftskip}{-2mm}
    \item Tìm loại có nhiều nhất, ít nhất.
    \item So sánh loại này nhiều hơn/ít hơn loại kia bao nhiêu.
    \item Tính tổng số đối tượng trên biểu đồ.
  \end{itemize}
\end{itemize}
}
\end{tcolorbox}
\end{minipage}

\vspace{-1.5mm}
\begin{tcolorbox}[
  enhanced,
  colback=orange!4!white,
  colframe=orange!80!black,
  arc=2mm, boxrule=0.8pt,
  left=2mm, right=2mm, top=1.5mm, bottom=1.5mm,
  title={\footnotesize\bfseries\color{white}\faSitemap\ 3. Mô phỏng cấu trúc trực quan của một biểu đồ tranh}
]
\centering
\begin{tikzpicture}[x=1.5cm, y=0.75cm]
  \tikzset{arr/.style={-{Stealth[scale=0.8]}}, every node/.append style={align=center}}
  \node[rectangle, rounded corners=2pt, draw=violet!80!black, fill=violet!15, font=\scriptsize\bfseries] at (2.5,0.7) {TÊN BIỂU ĐỒ: SỐ QUẢ TRÁI CÂY TRONG RỔ};

  \node[font=\tiny\bfseries, align=right] at (0.5,0) {Táo:};
  \node[font=\normalsize\bfseries, color=red!80!black] at (2.2,0) {$\bigstar\ \bigstar\ \bigstar\ \bigstar$};
  \node[font=\tiny\bfseries] at (4.2,0) {($4$ quả)};

  \node[font=\tiny\bfseries, align=right] at (0.5,-0.6) {Cam:};
  \node[font=\normalsize\bfseries, color=orange!90!black] at (2.5,-0.6) {$\bigstar\ \bigstar\ \bigstar\ \bigstar\ \bigstar$};
  \node[font=\tiny\bfseries] at (4.2,-0.6) {($5$ quả)};

  \node[font=\tiny\itshape, color=black!70] at (2.5,-1.2) {(Ghi chú: Mỗi $\bigstar$ biểu diễn $1$ quả)};
\end{tikzpicture}
\end{tcolorbox}

\vspace{-1.5mm}
\noindent
\begin{minipage}[t]{0.49\textwidth}\vspace{0pt}
\begin{tcolorbox}[
  enhanced,
  colback=green!4!white,
  colframe=green!60!black,
  arc=2mm, boxrule=0.8pt,
  left=2mm, right=2mm, top=1.8mm, bottom=1.8mm,
  title={\footnotesize\bfseries\color{white}\faCheckCircle\ 4. Câu hỏi khai thác biểu đồ mẫu}
]
{\footnotesize
Dựa vào biểu đồ trên, trả lời:\\
- Trong rổ có \textbf{$4$ quả táo} và \textbf{$5$ quả cam}.\\
- Loại quả nhiều hơn là: \textbf{Cam} (nhiều hơn $1$ quả vì $5 - 4 = 1$).\\
- Cả rổ có tất cả bao nhiêu quả?\\
$4 + 5 = \mathbf{9}$ (quả).\\
- Nếu ăn mất $2$ quả cam thì còn: $5 - 2 = \mathbf{3}$ quả cam.
}
\end{tcolorbox}
\end{minipage}
\hfill
\begin{minipage}[t]{0.49\textwidth}\vspace{0pt}
\begin{tcolorbox}[
  enhanced,
  colback=rose!4!white,
  colframe=rose!70!black,
  arc=2mm, boxrule=0.8pt,
  left=2mm, right=2mm, top=1.8mm, bottom=1.8mm,
  title={\footnotesize\bfseries\color{white}\faExclamationTriangle\ 5. Cạm bẫy học sinh cần nhớ}
]
{\footnotesize
\begin{itemize}\setlength{\itemsep}{1.5pt}\setlength{\leftskip}{-2.5mm}
  \item \textbf{Không đọc phần quy ước ghi chú:} Luôn phải nhìn chú thích ở dưới cùng xem $1$ hình vẽ ứng với mấy đồ vật.
  \item Nhầm lẫn giữa số biểu tượng tranh và số lượng thực tế khi có bài toán quy ước đặc biệt.
  \item Trả lời cộc lốc thiếu đơn vị (phải ghi rõ $4$ \textbf{quả táo}, $5$ \textbf{quả cam}).
\end{itemize}
}
\end{tcolorbox}
\end{minipage}


\end{document}
